\documentclass[11pt]{article}
\usepackage[margin=1in]{geometry}
\usepackage{amsmath,amssymb,amsthm}
\usepackage{xurl}
\usepackage{hyperref}
\sloppy
\let\oldus\_ \renewcommand{\_}{\oldus\allowbreak}
\newtheorem{theorem}{Theorem}
\newtheorem{lemma}[theorem]{Lemma}
\theoremstyle{definition}
\newtheorem{definition}[theorem]{Definition}
\newcommand{\tr}{\operatorname{tr}}

\title{Disproof of conjecture 00000000489:\\ $S_3(n)=\sum_{\lambda\vdash n} f_\lambda^3/n!$ is at most $\sqrt{n!}$,\\ so it is not $\Theta(\sqrt{n!}\,n^{1/4})$}
\author{Jackmeson1}
\date{2026-10-05}

\begin{document}
\maketitle

\section{The conjecture and the reading}

\textbf{Statement (English, verbatim).} Definition: The third-order character sum
$S_3(n)=\sum_\lambda f_\lambda^3/n!$. Conjecture: $S_3(n)=\Theta(\sqrt{n!}\cdot n^{1/4})$, and the
growth of $S_3(n)/\sqrt{n!}$ is given by the main term of the 3-thread counting of the limit shape,
with main-term constant $2^{1/4}\pi^{-1/2}$.

The Chinese text gives the same definition and, after a garbled prefix
(``$S_3(n)\sim c\cdot n^{1/4}\cdot(27/4)^{n^2/4\cdot 0+\dots}$''), says ``precisely
$S_3(n)=\Theta(\sqrt{n!}\cdot n^{1/4})$'' with the same main-term constant $2^{1/4}\pi^{-1/2}$.

\textbf{Objects.} $\lambda$ runs over the partitions of $n$, and $f_\lambda$ is the dimension of
the irreducible complex representation of the symmetric group $S_n$ labelled by $\lambda$. The
irreducible complex representations of $S_n$ are, up to isomorphism, in bijection with the
partitions of $n$ (Section~\ref{sec:sources}). Write $\sqrt{n!}=(n!)^{1/2}$, as in the
statement's own ``$S_3(n)/\sqrt{(n!)}$'', and let $n^{1/4}$ be the real positive fourth root.

\textbf{Claims refuted.} With $g(n)=\sqrt{n!}\,n^{1/4}$ and $C=2^{1/4}/\sqrt{\pi}$:
\begin{enumerate}
\item[(T)] $S_3(n)=\Theta(g(n))$ as $n\to\infty$, i.e.\ $S_3=O(g)$ and $g=O(S_3)$;
\item[(L)] $S_3(n)/g(n)\to C$;
\item[(E)] $S_3(n)/\sqrt{n!}\sim C\,n^{1/4}$.
\end{enumerate}
(L) and (E) are our readings of ``the growth of $S_3(n)/\sqrt{n!}$ \dots{} with main-term
constant $2^{1/4}\pi^{-1/2}$''. We do not give a meaning to ``3-thread counting of the limit
shape'' and we do not address the garbled Chinese prefix beyond its explicit $\Theta$ claim.
We prove $S_3(n)\le\sqrt{n!}$ for every $n$, hence $S_3(n)/g(n)\le n^{-1/4}\to 0$, which refutes
the lower half of (T), and (L) and (E).

\section{The mathematical proof}

Let $G$ be a finite group and $V_1,\dots,V_k$ pairwise non-isomorphic irreducible
finite-dimensional complex representations with characters $\chi_i$ and dimensions $d_i$.

\begin{lemma}[$\chi(g^{-1})=\overline{\chi(g)}$]\label{lem:conj}
Let $A:G\to M_d(\mathbb C)$ satisfy $A(gh)=A(g)A(h)$ and $A(1)=I$. Then
$\tr A(g^{-1})=\overline{\tr A(g)}$.
\end{lemma}
\begin{proof}
Put $P=\sum_{h\in G}A(h)^*A(h)$. Then $P=I+\sum_{h\ne1}A(h)^*A(h)$ is positive definite, hence
invertible. Reindexing $h\mapsto hg$ gives $A(g)^*PA(g)=P$. Since $A(g)A(g^{-1})=I$, this gives
$A(g)^*P=PA(g^{-1})$, so $A(g)^*=PA(g^{-1})P^{-1}$ and
$\overline{\tr A(g)}=\tr A(g)^*=\tr A(g^{-1})$.
\end{proof}
Applying the lemma to the matrices of $V$ in a basis gives $\chi_V(g^{-1})=\overline{\chi_V(g)}$.

\begin{lemma}\label{lem:bessel}
$\sum_i d_i^2\le|G|$.
\end{lemma}
\begin{proof}
By orthogonality of characters,
$\sum_{g}\chi_i(g)\chi_j(g^{-1})=|G|\,[i=j]$, because $V_i\cong V_j$ iff $i=j$. Let
$\Phi=\sum_i d_i\chi_i$ and $D=\sum_i d_i^2$. Then $\sum_g\Phi(g)\Phi(g^{-1})=|G|D$. By
Lemma~\ref{lem:conj}, $\Phi(g^{-1})=\overline{\Phi(g)}$, so
$|G|D=\sum_g|\Phi(g)|^2\ge|\Phi(1)|^2=D^2$, using $\chi_i(1)=d_i$. Hence $D\le|G|$.
\end{proof}

\begin{theorem}\label{thm:main}
$\sum_i d_i^3\le\sqrt{|G|}\,|G|$. For $G=S_n$ and $d_i=f_\lambda$ this is $S_3(n)\le\sqrt{n!}$.
\end{theorem}
\begin{proof}
Each $d_i^2\le D\le|G|$, so $d_i\le\sqrt{|G|}$ and $\sum_i d_i^3\le\sqrt{|G|}\sum_i d_i^2\le
\sqrt{|G|}\,|G|$. Divide by $|G|=n!$.
\end{proof}

Hence $0\le S_3(n)/g(n)\le n^{-1/4}\to0$, so $S_3=o(g)$. If $g=O(S_3)$ held, then $g=o(g)$, which
is impossible since $g(n)\ne0$ for $n\ge1$. So (T) fails. (L) fails because the limit is
$0\ne C$. For (E): $\bigl(S_3(n)/\sqrt{n!}\bigr)/\bigl(C n^{1/4}\bigr)=\bigl(S_3(n)/g(n)\bigr)/C\to
0\ne1$.

\section{Lean formalization}

The file \texttt{lean/Conjecture489/Basic.lean} (264 lines, \texttt{import Mathlib} only, Lean
v4.33.1, Mathlib v4.33.1) proves the following in namespace \texttt{C489}.
\begin{itemize}
\item \texttt{trace\_inv\_eq\_star} and \texttt{char\_inv}: Lemma~\ref{lem:conj} and
  $\chi_V(g^{-1})=\overline{\chi_V(g)}$ for every \texttt{V : FDRep $\mathbb C$ G}.
\item \texttt{sum\_sq\_finrank\_le}: for a finite type $\iota$ and
  \texttt{V : $\iota$ $\to$ FDRep $\mathbb C$ G} with every \texttt{V i} \texttt{Simple} and
  \texttt{Nonempty (V i $\cong$ V j) $\to$ i = j}, $\sum_i(\dim V_i)^2\le|G|$. It uses Mathlib's
  \texttt{FDRep.char\_orthonormal} and \texttt{FDRep.char\_one}.
\item \texttt{sum\_cube\_le}: $\sum_i(\dim V_i)^3\le\sqrt{|G|}\,|G|$.
\item \texttt{IrrepFamily n}: a map \texttt{V : Nat.Partition n $\to$ FDRep $\mathbb C$
  (Equiv.Perm (Fin n))}, with each $V_\lambda$ simple and distinct labels non-isomorphic.
  \texttt{IrrepFamily.f} is $f_\lambda=\dim_{\mathbb C}V_\lambda$ (\texttt{finrank}), and
  \texttt{IrrepFamily.S3} is $\bigl(\sum_\lambda f_\lambda^3\bigr)/n!$.
\item \texttt{S3\_le\_sqrt}: $S_3(n)\le\sqrt{n!}$ for every \texttt{IrrepFamily n}.
\item \texttt{conjecture489\_false}: for every \texttt{F : $\forall$ n, IrrepFamily n},
  (i) $S_3(n)\le\sqrt{n!}$ for all $n$; (ii) $S_3(n)/(\sqrt{n!}\,n^{1/4})\to0$;
  (iii) $S_3=o(g)$ (\texttt{=o[atTop]});
  (iv) $\neg\,S_3=\Theta(g)$ (\texttt{=$\Theta$[atTop]}); (v) the ratio does not tend to
  $2^{1/4}/\sqrt\pi$; (vi) $\neg\,\bigl(S_3(n)/\sqrt{n!}\bigr)\sim C n^{1/4}$
  (\texttt{\textasciitilde[atTop]}).
\end{itemize}
The theorem holds for every labelled family of pairwise non-isomorphic irreducibles. It does not
assume that the family is complete, and it does not construct Specht modules. The family that
defines the conjecture's $f_\lambda$ (one irreducible for each partition, pairwise
non-isomorphic) is one such family, so the conjecture's $S_3(n)$ is covered. Here
$n^{1/4}$ is \texttt{Real.rpow}, and $0^{1/4}=0$, so $g(0)=0$; only $n\ge1$ matters for the
asymptotic claims.

\textbf{Axiom audit.} \texttt{lake build} succeeds, and \texttt{\#print axioms} for
\texttt{C489.conjecture489\_false}, \texttt{C489.sum\_sq\_finrank\_le} and
\texttt{C489.char\_inv} prints \texttt{[propext, Classical.choice, Quot.sound]}. The file contains no
\texttt{sorry}, no \texttt{native\_decide} and no added axioms.

\section{Sources}\label{sec:sources}
Wikipedia, ``Representation theory of the symmetric group''
(\url{https://en.wikipedia.org/wiki/Representation_theory_of_the_symmetric_group}), retrieved
2026-10-05:
``Therefore according to the representation theory of a finite group, the number of inequivalent
irreducible representations, over the complex numbers, is equal to the number of partitions of n.
Unlike the general situation for finite groups, there is in fact a natural way to parametrize
irreducible representations by the same set that parametrizes conjugacy classes, namely by
partitions of n or equivalently Young diagrams of size n.''
The same article writes $d_\lambda$ for ``the dimension \dots{} of the representation that
corresponds to the Young diagram $\lambda$''. No other source is used. All other steps are proved
above and in Lean.

\end{document}
